% Options for packages loaded elsewhere
% Options for packages loaded elsewhere
\PassOptionsToPackage{unicode}{hyperref}
\PassOptionsToPackage{hyphens}{url}
\PassOptionsToPackage{dvipsnames,svgnames,x11names}{xcolor}
%
\documentclass[
  letterpaper,
  DIV=11,
  numbers=noendperiod]{scrartcl}
\usepackage{xcolor}
\usepackage{amsmath,amssymb}
\setcounter{secnumdepth}{5}
\usepackage{iftex}
\ifPDFTeX
  \usepackage[T1]{fontenc}
  \usepackage[utf8]{inputenc}
  \usepackage{textcomp} % provide euro and other symbols
\else % if luatex or xetex
  \usepackage{unicode-math} % this also loads fontspec
  \defaultfontfeatures{Scale=MatchLowercase}
  \defaultfontfeatures[\rmfamily]{Ligatures=TeX,Scale=1}
\fi
\usepackage{lmodern}
\ifPDFTeX\else
  % xetex/luatex font selection
\fi
% Use upquote if available, for straight quotes in verbatim environments
\IfFileExists{upquote.sty}{\usepackage{upquote}}{}
\IfFileExists{microtype.sty}{% use microtype if available
  \usepackage[]{microtype}
  \UseMicrotypeSet[protrusion]{basicmath} % disable protrusion for tt fonts
}{}
\makeatletter
\@ifundefined{KOMAClassName}{% if non-KOMA class
  \IfFileExists{parskip.sty}{%
    \usepackage{parskip}
  }{% else
    \setlength{\parindent}{0pt}
    \setlength{\parskip}{6pt plus 2pt minus 1pt}}
}{% if KOMA class
  \KOMAoptions{parskip=half}}
\makeatother
% Make \paragraph and \subparagraph free-standing
\makeatletter
\ifx\paragraph\undefined\else
  \let\oldparagraph\paragraph
  \renewcommand{\paragraph}{
    \@ifstar
      \xxxParagraphStar
      \xxxParagraphNoStar
  }
  \newcommand{\xxxParagraphStar}[1]{\oldparagraph*{#1}\mbox{}}
  \newcommand{\xxxParagraphNoStar}[1]{\oldparagraph{#1}\mbox{}}
\fi
\ifx\subparagraph\undefined\else
  \let\oldsubparagraph\subparagraph
  \renewcommand{\subparagraph}{
    \@ifstar
      \xxxSubParagraphStar
      \xxxSubParagraphNoStar
  }
  \newcommand{\xxxSubParagraphStar}[1]{\oldsubparagraph*{#1}\mbox{}}
  \newcommand{\xxxSubParagraphNoStar}[1]{\oldsubparagraph{#1}\mbox{}}
\fi
\makeatother

\usepackage{color}
\usepackage{fancyvrb}
\newcommand{\VerbBar}{|}
\newcommand{\VERB}{\Verb[commandchars=\\\{\}]}
\DefineVerbatimEnvironment{Highlighting}{Verbatim}{commandchars=\\\{\}}
% Add ',fontsize=\small' for more characters per line
\usepackage{framed}
\definecolor{shadecolor}{RGB}{241,243,245}
\newenvironment{Shaded}{\begin{snugshade}}{\end{snugshade}}
\newcommand{\AlertTok}[1]{\textcolor[rgb]{0.68,0.00,0.00}{#1}}
\newcommand{\AnnotationTok}[1]{\textcolor[rgb]{0.37,0.37,0.37}{#1}}
\newcommand{\AttributeTok}[1]{\textcolor[rgb]{0.40,0.45,0.13}{#1}}
\newcommand{\BaseNTok}[1]{\textcolor[rgb]{0.68,0.00,0.00}{#1}}
\newcommand{\BuiltInTok}[1]{\textcolor[rgb]{0.00,0.23,0.31}{#1}}
\newcommand{\CharTok}[1]{\textcolor[rgb]{0.13,0.47,0.30}{#1}}
\newcommand{\CommentTok}[1]{\textcolor[rgb]{0.37,0.37,0.37}{#1}}
\newcommand{\CommentVarTok}[1]{\textcolor[rgb]{0.37,0.37,0.37}{\textit{#1}}}
\newcommand{\ConstantTok}[1]{\textcolor[rgb]{0.56,0.35,0.01}{#1}}
\newcommand{\ControlFlowTok}[1]{\textcolor[rgb]{0.00,0.23,0.31}{\textbf{#1}}}
\newcommand{\DataTypeTok}[1]{\textcolor[rgb]{0.68,0.00,0.00}{#1}}
\newcommand{\DecValTok}[1]{\textcolor[rgb]{0.68,0.00,0.00}{#1}}
\newcommand{\DocumentationTok}[1]{\textcolor[rgb]{0.37,0.37,0.37}{\textit{#1}}}
\newcommand{\ErrorTok}[1]{\textcolor[rgb]{0.68,0.00,0.00}{#1}}
\newcommand{\ExtensionTok}[1]{\textcolor[rgb]{0.00,0.23,0.31}{#1}}
\newcommand{\FloatTok}[1]{\textcolor[rgb]{0.68,0.00,0.00}{#1}}
\newcommand{\FunctionTok}[1]{\textcolor[rgb]{0.28,0.35,0.67}{#1}}
\newcommand{\ImportTok}[1]{\textcolor[rgb]{0.00,0.46,0.62}{#1}}
\newcommand{\InformationTok}[1]{\textcolor[rgb]{0.37,0.37,0.37}{#1}}
\newcommand{\KeywordTok}[1]{\textcolor[rgb]{0.00,0.23,0.31}{\textbf{#1}}}
\newcommand{\NormalTok}[1]{\textcolor[rgb]{0.00,0.23,0.31}{#1}}
\newcommand{\OperatorTok}[1]{\textcolor[rgb]{0.37,0.37,0.37}{#1}}
\newcommand{\OtherTok}[1]{\textcolor[rgb]{0.00,0.23,0.31}{#1}}
\newcommand{\PreprocessorTok}[1]{\textcolor[rgb]{0.68,0.00,0.00}{#1}}
\newcommand{\RegionMarkerTok}[1]{\textcolor[rgb]{0.00,0.23,0.31}{#1}}
\newcommand{\SpecialCharTok}[1]{\textcolor[rgb]{0.37,0.37,0.37}{#1}}
\newcommand{\SpecialStringTok}[1]{\textcolor[rgb]{0.13,0.47,0.30}{#1}}
\newcommand{\StringTok}[1]{\textcolor[rgb]{0.13,0.47,0.30}{#1}}
\newcommand{\VariableTok}[1]{\textcolor[rgb]{0.07,0.07,0.07}{#1}}
\newcommand{\VerbatimStringTok}[1]{\textcolor[rgb]{0.13,0.47,0.30}{#1}}
\newcommand{\WarningTok}[1]{\textcolor[rgb]{0.37,0.37,0.37}{\textit{#1}}}

\usepackage{longtable,booktabs,array}
\usepackage{calc} % for calculating minipage widths
% Correct order of tables after \paragraph or \subparagraph
\usepackage{etoolbox}
\makeatletter
\patchcmd\longtable{\par}{\if@noskipsec\mbox{}\fi\par}{}{}
\makeatother
% Allow footnotes in longtable head/foot
\IfFileExists{footnotehyper.sty}{\usepackage{footnotehyper}}{\usepackage{footnote}}
\makesavenoteenv{longtable}
\usepackage{graphicx}
\makeatletter
\newsavebox\pandoc@box
\newcommand*\pandocbounded[1]{% scales image to fit in text height/width
  \sbox\pandoc@box{#1}%
  \Gscale@div\@tempa{\textheight}{\dimexpr\ht\pandoc@box+\dp\pandoc@box\relax}%
  \Gscale@div\@tempb{\linewidth}{\wd\pandoc@box}%
  \ifdim\@tempb\p@<\@tempa\p@\let\@tempa\@tempb\fi% select the smaller of both
  \ifdim\@tempa\p@<\p@\scalebox{\@tempa}{\usebox\pandoc@box}%
  \else\usebox{\pandoc@box}%
  \fi%
}
% Set default figure placement to htbp
\def\fps@figure{htbp}
\makeatother





\setlength{\emergencystretch}{3em} % prevent overfull lines

\providecommand{\tightlist}{%
  \setlength{\itemsep}{0pt}\setlength{\parskip}{0pt}}



 


\KOMAoption{captions}{tableheading}
\makeatletter
\@ifpackageloaded{caption}{}{\usepackage{caption}}
\AtBeginDocument{%
\ifdefined\contentsname
  \renewcommand*\contentsname{Table of contents}
\else
  \newcommand\contentsname{Table of contents}
\fi
\ifdefined\listfigurename
  \renewcommand*\listfigurename{List of Figures}
\else
  \newcommand\listfigurename{List of Figures}
\fi
\ifdefined\listtablename
  \renewcommand*\listtablename{List of Tables}
\else
  \newcommand\listtablename{List of Tables}
\fi
\ifdefined\figurename
  \renewcommand*\figurename{Figure}
\else
  \newcommand\figurename{Figure}
\fi
\ifdefined\tablename
  \renewcommand*\tablename{Table}
\else
  \newcommand\tablename{Table}
\fi
}
\@ifpackageloaded{float}{}{\usepackage{float}}
\floatstyle{ruled}
\@ifundefined{c@chapter}{\newfloat{codelisting}{h}{lop}}{\newfloat{codelisting}{h}{lop}[chapter]}
\floatname{codelisting}{Listing}
\newcommand*\listoflistings{\listof{codelisting}{List of Listings}}
\makeatother
\makeatletter
\makeatother
\makeatletter
\@ifpackageloaded{caption}{}{\usepackage{caption}}
\@ifpackageloaded{subcaption}{}{\usepackage{subcaption}}
\makeatother
\usepackage{bookmark}
\IfFileExists{xurl.sty}{\usepackage{xurl}}{} % add URL line breaks if available
\urlstyle{same}
\hypersetup{
  pdftitle={A Small Open-Economy DSGE Model for the UK},
  pdfauthor={Working Paper No.~1},
  colorlinks=true,
  linkcolor={blue},
  filecolor={Maroon},
  citecolor={Blue},
  urlcolor={Blue},
  pdfcreator={LaTeX via pandoc}}


\title{A Small Open-Economy DSGE Model for the UK}
\usepackage{etoolbox}
\makeatletter
\providecommand{\subtitle}[1]{% add subtitle to \maketitle
  \apptocmd{\@title}{\par {\large #1 \par}}{}{}
}
\makeatother
\subtitle{Stage 1: System Estimation and Preliminary Forecasting}
\author{Working Paper No.~1}
\date{2026-03-14}
\begin{document}
\maketitle

\renewcommand*\contentsname{Table of contents}
{
\hypersetup{linkcolor=}
\setcounter{tocdepth}{3}
\tableofcontents
}

\section{1. Introduction}\label{introduction}

\ldots{}

\section{2. The Theoretical Structural
Model}\label{the-theoretical-structural-model}

Following the structural modeling tradition of @Scarth1988, we define
the UK economy as a Small Open Economy (SOE) characterized by nominal
rigidities and rational expectations. The model consists of a system of
stochastic behavioural equations and non-stochastic national accounting
identities.

\subsection{2.1 Stochastic Behavioural
Equations}\label{stochastic-behavioural-equations}

The core dynamics are driven by forward-looking agents and the
Mundell-Fleming extension for open-economy linkages. All variables
(unless otherwise noted) are expressed as log-deviations from their
steady-state trends.

\textbf{The Open-Economy IS Curve:} Aggregate demand depends on expected
future output, real interest rates, and the real exchange rate (the
expenditure-switching effect):
\[y_t = E_t y_{t+1} - \sigma(r_t - E_t \pi_{t+1}) + \gamma q_t + \epsilon^y_t\]
where \(q_t\) is the real exchange rate and \(\sigma\) represents the
intertemporal elasticity of substitution.

\textbf{The Rational Expectations Phillips Curve (REPC):} Inflation is
driven by forward-looking expectations and domestic demand pressure (the
output gap):
\[\pi_t = \beta E_t \pi_{t+1} + \kappa y_t + \epsilon^\pi_t\]

\textbf{Uncovered Interest Parity (UIP):} To close the capital account
under rational expectations, the domestic interest rate differential
must equal the expected depreciation of the currency:
\[r_t - r^*_t = E_t e_{t+1} - e_t + \phi_t\] where \(\phi_t\) represents
a stochastic risk premium.

\subsection{2.2 National Accounting
Identities}\label{national-accounting-identities}

A critical departure in this paper is the explicit inclusion of national
accounting identities as non-stochastic constraints within the 3SLS
estimation framework. These identities ensure the structural model
remains consistent with the ONS dataset.

\textbf{The Production-Tax Identity:} We define Gross Domestic Product
(\(y_t\)) as the sum of the Gross Value Added of industries (\(y^d_t\))
and net taxes on products (\(T_t\)):
\[y_t = \omega y^d_t + (1-\omega) T_t\] where \(\omega\) represents the
steady-state share of GVA in total output.

\textbf{The Expenditure Identity:} To ensure the IS story is fully
closed, aggregate output must also satisfy the expenditure equilibrium:
\[y_t = c_{y}c_t + i_{y}i_t + g_{y}g_t + nx_t\] where the coefficients
denote the shares of consumption, investment, government spending, and
net exports in total GDP.

By defining both the production side (\(y^d + T\)) and the expenditure
side (\(c + i + g + nx\)), the model accounts for the ``double-entry''
nature of the UK national accounts, providing a more robust set of
instruments for the 3SLS estimation.

\section{3. The empirical model}\label{the-empirical-model}

\subsection{3.1 Variables}\label{variables}

check against formulation above, but this is a working list of variables

\begin{Shaded}
\begin{Highlighting}[]
\FunctionTok{library}\NormalTok{(knitr)}
\FunctionTok{library}\NormalTok{(tibble)}

\CommentTok{\# We use double backslashes (\textbackslash{}\textbackslash{}) so R treats them as literal LaTeX commands}
\NormalTok{var\_mapping }\OtherTok{\textless{}{-}} \FunctionTok{tribble}\NormalTok{(}
  \SpecialCharTok{\textasciitilde{}}\StringTok{"Math Notation"}\NormalTok{, }\SpecialCharTok{\textasciitilde{}}\StringTok{"R Variable"}\NormalTok{,     }\SpecialCharTok{\textasciitilde{}}\StringTok{"Description"}\NormalTok{,                }\SpecialCharTok{\textasciitilde{}}\StringTok{"Classification"}\NormalTok{,}
  \StringTok{"$y\_t$"}\NormalTok{,          }\StringTok{"Y\_GAP"}\NormalTok{,          }\StringTok{"Output Gap"}\NormalTok{,                  }\StringTok{"Endogenous"}\NormalTok{,}
  \StringTok{"$}\SpecialCharTok{\textbackslash{}\textbackslash{}}\StringTok{pi\_t$"}\NormalTok{,       }\StringTok{"Infl\_dev"}\NormalTok{,       }\StringTok{"Inflation Deviation"}\NormalTok{,         }\StringTok{"Endogenous"}\NormalTok{,}
  \StringTok{"$}\SpecialCharTok{\textbackslash{}\textbackslash{}}\StringTok{Delta i\_t$"}\NormalTok{,  }\StringTok{"d\_BOE\_rate"}\NormalTok{,     }\StringTok{"Change in Bank Rate"}\NormalTok{,         }\StringTok{"Endogenous"}\NormalTok{,}
  \StringTok{"$u\_t$"}\NormalTok{,          }\StringTok{"U\_GAP"}\NormalTok{,          }\StringTok{"Unemployment Gap"}\NormalTok{,            }\StringTok{"Endogenous"}\NormalTok{,}
  \StringTok{"$r\_t$"}\NormalTok{,          }\StringTok{"real\_BOE\_r"}\NormalTok{,     }\StringTok{"Ex{-}ante Real Interest Rate"}\NormalTok{,  }\StringTok{"Exogenous"}\NormalTok{,}
  \StringTok{"$spr\_t$"}\NormalTok{,        }\StringTok{"d\_mkt\_spr"}\NormalTok{,      }\StringTok{"Market Credit Spread"}\NormalTok{,        }\StringTok{"Exogenous"}\NormalTok{,}
  \StringTok{"$}\SpecialCharTok{\textbackslash{}\textbackslash{}}\StringTok{Delta nx\_t$"}\NormalTok{, }\StringTok{"dlog\_NX"}\NormalTok{,        }\StringTok{"Growth in Net Exports"}\NormalTok{,       }\StringTok{"Exogenous"}\NormalTok{,}
  \StringTok{"$q\_t$"}\NormalTok{,          }\StringTok{"rer\_gap"}\NormalTok{,        }\StringTok{"Real Exchange Rate Gap"}\NormalTok{,      }\StringTok{"Exogenous"}\NormalTok{,}
  \StringTok{"$}\SpecialCharTok{\textbackslash{}\textbackslash{}}\StringTok{pi\^{}\{e\}\_t$"}\NormalTok{,   }\StringTok{"dlog\_energy\_cpi"}\NormalTok{,}\StringTok{"Energy Price Inflation"}\NormalTok{,      }\StringTok{"Instrument"}\NormalTok{,}
  \StringTok{"$}\SpecialCharTok{\textbackslash{}\textbackslash{}}\StringTok{Delta m\_t$"}\NormalTok{,  }\StringTok{"dlog\_IM\_defl"}\NormalTok{,   }\StringTok{"Import Price Growth"}\NormalTok{,         }\StringTok{"Instrument"}\NormalTok{,}
  \StringTok{"$}\SpecialCharTok{\textbackslash{}\textbackslash{}}\StringTok{Delta c\_t$"}\NormalTok{,  }\StringTok{"dlog\_cap"}\NormalTok{,       }\StringTok{"Capacity Utilization Growth"}\NormalTok{, }\StringTok{"Instrument"}\NormalTok{,}
  \StringTok{"$D\_\{2020\}$"}\NormalTok{,     }\StringTok{"dummy\_2020Q2/3"}\NormalTok{, }\StringTok{"Structural Dummies (COVID)"}\NormalTok{,  }\StringTok{"Exogenous"}
\NormalTok{)}

\FunctionTok{kable}\NormalTok{(var\_mapping, }\AttributeTok{escape =} \ConstantTok{FALSE}\NormalTok{)}
\end{Highlighting}
\end{Shaded}

\begin{longtable}[]{@{}
  >{\raggedright\arraybackslash}p{(\linewidth - 6\tabcolsep) * \real{0.1918}}
  >{\raggedright\arraybackslash}p{(\linewidth - 6\tabcolsep) * \real{0.2192}}
  >{\raggedright\arraybackslash}p{(\linewidth - 6\tabcolsep) * \real{0.3836}}
  >{\raggedright\arraybackslash}p{(\linewidth - 6\tabcolsep) * \real{0.2055}}@{}}

\caption{\label{tbl-variable-mapping}System Variables: Mapping and
Classification for Working Paper No.~1}

\tabularnewline

\toprule\noalign{}
\begin{minipage}[b]{\linewidth}\raggedright
Math Notation
\end{minipage} & \begin{minipage}[b]{\linewidth}\raggedright
R Variable
\end{minipage} & \begin{minipage}[b]{\linewidth}\raggedright
Description
\end{minipage} & \begin{minipage}[b]{\linewidth}\raggedright
Classification
\end{minipage} \\
\midrule\noalign{}
\endhead
\bottomrule\noalign{}
\endlastfoot
\(y_t\) & Y\_GAP & Output Gap & Endogenous \\
\(\pi_t\) & Infl\_dev & Inflation Deviation & Endogenous \\
\(\Delta i_t\) & d\_BOE\_rate & Change in Bank Rate & Endogenous \\
\(u_t\) & U\_GAP & Unemployment Gap & Endogenous \\
\(r_t\) & real\_BOE\_r & Ex-ante Real Interest Rate & Exogenous \\
\(spr_t\) & d\_mkt\_spr & Market Credit Spread & Exogenous \\
\(\Delta nx_t\) & dlog\_NX & Growth in Net Exports & Exogenous \\
\(q_t\) & rer\_gap & Real Exchange Rate Gap & Exogenous \\
\(\pi^{e}_t\) & dlog\_energy\_cpi & Energy Price Inflation &
Instrument \\
\(\Delta m_t\) & dlog\_IM\_defl & Import Price Growth & Instrument \\
\(\Delta c_t\) & dlog\_cap & Capacity Utilization Growth & Instrument \\
\(D_{2020}\) & dummy\_2020Q2/3 & Structural Dummies (COVID) &
Exogenous \\

\end{longtable}

\subsection{3.2 Structural Equations}\label{structural-equations}

The structural system for our model is defined as follows:

\subsubsection{National Accounting
Constraints}\label{national-accounting-constraints}

To ensure the empirical model is consistent with UK National Accounts,
the estimated output gap (\(y_t\)) is subject to two fundamental
identities:

\[
y_t = \omega y^d_t + (1-\omega) \tau_t + \xi_{p,t} \tag{Production Constraint}
\]

\[
y_t = w_c c_t + w_i i_t + w_g g_t + w_{nx} nx_t + \xi_{e,t} \tag{Expenditure Constraint}
\]

Where \(\omega\) and \(w_j\) are the steady-state shares derived from
the ONS dataset, and \(\xi_t\) represents the statistical discrepancy
(the `check' variables in our data).

\subsubsection{Behavioral Equations}\label{behavioral-equations}

The core dynamics are captured by the following four-equation system,
which incorporates both forward-looking rational expectations and
backward-looking persistence:

\[
\begin{aligned}
\text{IS Curve:} \quad & y_t = \beta_1 y_{t-1} + \beta_2 y_{t+1|t} + \beta_3 r_t + \beta_4 spr_t + \beta_5 \Delta nx_t + \beta_6 q_{t-1} + \mathbf{D}\gamma_{IS} + \epsilon_{y,t} \\
\text{Phillips Curve:} \quad & \pi_t = \delta_1 \pi_{t-1} + \delta_2 \pi_{t-2} + \delta_3 y_t + \delta_4 \pi^{e}_t + \delta_5 \Delta m_t + \delta_6 \Delta c_t + \epsilon_{\pi,t} \\
\text{Taylor Rule:} \quad & \Delta i_t = \theta_1 i_{t-1} + \theta_2 \pi_t + \theta_3 y_t + \theta_4 spr_t + \gamma_{3} D_{20Q2} + \epsilon_{i,t} \\
\text{Okun's Law:} \quad & u_t = \alpha_1 u_{t-1} + \alpha_2 y_t + \epsilon_{u,t}
\end{aligned}
\]

Where \(\omega\) and \(w_j\) are the steady-state shares derived from
the ONS dataset, and \(\xi_t\) represents the statistical discrepancy
(the `check' variables in our data).

\subsection{3.3 Data}\label{data}

\begin{Shaded}
\begin{Highlighting}[]
\NormalTok{DSGE\_3SLS\_DATA }\OtherTok{\textless{}{-}} \FunctionTok{read.csv}\NormalTok{(}\StringTok{"C:}\SpecialCharTok{\textbackslash{}\textbackslash{}}\StringTok{Users}\SpecialCharTok{\textbackslash{}\textbackslash{}}\StringTok{hogan}\SpecialCharTok{\textbackslash{}\textbackslash{}}\StringTok{OneDrive}\SpecialCharTok{\textbackslash{}\textbackslash{}}\StringTok{Documents}\SpecialCharTok{\textbackslash{}\textbackslash{}}\StringTok{1 OLLIE}\SpecialCharTok{\textbackslash{}\textbackslash{}}\StringTok{0. R}\SpecialCharTok{\textbackslash{}\textbackslash{}}\StringTok{UK\_economic\_model}\SpecialCharTok{\textbackslash{}\textbackslash{}}\StringTok{reports}\SpecialCharTok{\textbackslash{}\textbackslash{}}\StringTok{MODEL\_DIFFS.csv"}\NormalTok{)}
\end{Highlighting}
\end{Shaded}

\subsection{3.4 3SLS model estimation}\label{sls-model-estimation}

\subsection{3.5 Review of robustness}\label{review-of-robustness}

\section{4. Simulations and impulse response
functions}\label{simulations-and-impulse-response-functions}

\ldots{}

\section{5. Forecasts and forecasting
performance}\label{forecasts-and-forecasting-performance}

\section{6. Conclusion and next steps}\label{conclusion-and-next-steps}




\end{document}
